\section{Scope of the Fresh V2 Results}

This chapter is the principal Fresh Full891 V2 result layer. It uses the 890 complete California architectures and, for central-node marginal curves, the 880 cases belonging to 80 complete eleven-level sweeps. The omitted case and all scenario counts are documented in Chapter~5 and Appendix~C. India Fresh V2 is not yet available; no California value in this chapter is presented as an India result.

Two scenarios anchor the interpretation. \emph{Nominal original} is the bounded central-only targeted policy with the operational deadline, three-hop limit, three observer copies, and five central-node copies. \emph{Unlimited useful deadline} removes the hop/copy limits while retaining the useful-recipient and deadline definitions. The second scenario is a diagnostic upper-capability replay, not a directly recommended operating mode.

\section{Bounded Policy versus Architecture Potential}

Table~\ref{tab:nominal_unlimited} quantifies the gap between the two anchor scenarios. Nominal mean observation fulfilment is already 61.59 percent, but strict useful-recipient completion is only 0.63 percent and mean useful coverage only 15.29 percent. Removing copy and hop limits raises strict completion to 74.69 percent and coverage to 91.14 percent. The increase requires more than ten times the transferred data.

\begin{table}[htbp]
\centering
\caption{California mean outcome over 890 complete Fresh V2 architectures.}
\label{tab:nominal_unlimited}
\begin{tabularx}{\textwidth}{L{0.28\textwidth}rrrr}
\toprule
Scenario & Obs. (\%) & \Suseful\ (\%) & Coverage (\%) & Data (Mbit/case) \\
\midrule
Nominal original & 61.59 & 0.63 & 15.29 & 197.55 \\
Unlimited useful deadline & 77.62 & 74.69 & 91.14 & 2,047.57 \\
\midrule
Difference & +16.04 & +74.07 & +75.85 & +1,850.02 \\
\bottomrule
\end{tabularx}
\end{table}

The main explanation is operational rather than purely geometric. The bounded policy frequently reaches a useful observer, which can complete the observation objective, but does not distribute the task to every useful recipient. The diagnostic replay uses additional copies and deeper temporal paths. Therefore, the low nominal \Suseful\ is a property of the architecture--policy combination, not evidence that the constellation has no communication opportunities.

\section{Central-Node Fraction}

Figure~\ref{fig:v2_cn_fraction} presents both operating regimes. Under the nominal policy, mean observation fulfilment increases sharply from 29.14 percent at CN0 to 64.88 percent at CN10, reaches 71.13 percent at CN50, and falls to 35.70 percent at CN100. Strict dissemination peaks at 1.72 percent at CN10, while mean useful coverage peaks around CN20 at approximately 20.7 percent and falls to 2.42 percent at CN100. Intermediate centralization is therefore preferable under the bounded policy.

\begin{figure}[htbp]
\centering
\includegraphics[width=\textwidth]{02_cn_fraction_performance.png}
\caption{Central-node fraction response across 80 balanced California Fresh V2 base architectures. Shaded or interval displays summarize variation across base geometries.}
\label{fig:v2_cn_fraction}
\end{figure}

The unlimited replay reveals a different, saturating regime. Selected marginal means are shown in Table~\ref{tab:v2_cn_selected}. \Suseful\ grows rapidly up to the intermediate range and then provides smaller incremental gains, while traffic continues to grow. From CN20 to CN100, mean \Suseful\ increases by about 14.6 points, but transferred data increases by approximately 2,429 Mbit per case (about 211 percent).

\begin{table}[htbp]
\centering
\caption{Selected unlimited useful-deadline central-node marginal means.}
\label{tab:v2_cn_selected}
\begin{tabular}{rrrrrr}
\toprule
CN (\%) & \Suseful\ (\%) & Coverage (\%) & P100 $T_{100}$ (min) & Data (Mbit) & Events \\
\midrule
0 & 5.16 & 45.21 & 232.6 & 297 & 407 \\
20 & 74.34 & 92.02 & 120.0 & 1,152 & 1,333 \\
40 & 83.40 & 96.61 & 97.4 & 1,771 & 1,939 \\
70 & 86.56 & -- & -- & -- & -- \\
100 & 88.95 & 98.55 & 73.2 & 3,581 & 3,709 \\
\bottomrule
\end{tabular}
\end{table}

The dashes denote values not selected for the compact table; they are not treated as zero. P100 $T_{100}$ is conditional and especially censored at CN0. In the nominal scenario, mean maximum routing depth reaches the configured limit of three by approximately CN20. In the unlimited replay it rises from 1.0 at CN0 to about 7.8 at CN20, 9.6 at CN40, and 11.8 at CN100. Part of the unlimited gain is therefore attributable to policy freedom, not central-node fraction alone.

\section{Architecture Factors and Dependence}

Figure~\ref{fig:v2_arch_effects} shows descriptive main effects. Within the unlimited balanced subset, unadjusted one-factor $\eta^2$ for CN fraction is 0.690 for \Suseful\ and 0.820 for useful coverage. Plane count is the next largest reported factor for \Suseful\ at 0.112. For observation fulfilment, CN fraction and plane count have $\eta^2$ values of 0.402 and 0.361. Satellite count, altitude, and inclination are smaller in these marginal summaries.

\begin{figure}[htbp]
\centering
\includegraphics[width=\textwidth]{03_architecture_main_effects.png}
\caption{California Fresh V2 marginal architecture effects under the unlimited useful-deadline replay.}
\label{fig:v2_arch_effects}
\end{figure}

These values show that dissemination depends strongly on the central-node role, while observation also depends materially on plane geometry. Because the design contains interactions and discrete placement changes, the unadjusted effects are descriptive. They do not imply that changing one factor while holding an unmodelled system fixed will reproduce the marginal mean.

\section{Policy Ablation}

The policy ablation in Table~\ref{tab:policy_ablation} separates central-only access from all-peer link availability. Peer-only and hybrid produce identical recorded outcomes in the current implementation. This identity is important: it indicates that the implemented hybrid preference does not change admitted transfers relative to peer-only for this scenario set. It is not evidence that every possible hybrid algorithm is equivalent to peer routing.

\begin{table}[htbp]
\centering
\caption{Mean California policy-ablation outcome over 890 complete architectures.}
\label{tab:policy_ablation}
\begin{tabular}{lrrrr}
\toprule
Policy scenario & Obs. (\%) & \Suseful\ (\%) & Coverage (\%) & Data (Mbit/case) \\
\midrule
Central / unlimited & 77.62 & 74.69 & 91.14 & 2,047.57 \\
Peer & 83.04 & 85.44 & 97.20 & 3,485.72 \\
Hybrid & 83.04 & 85.44 & 97.20 & 3,485.72 \\
\bottomrule
\end{tabular}
\end{table}

\begin{figure}[htbp]
\centering
\includegraphics[width=0.92\textwidth]{04_policy_ablation.png}
\caption{Policy ablation: performance gain and communication burden under central, peer, and hybrid link use.}
\label{fig:policy_ablation}
\end{figure}

The peer-capable variants improve observation by 5.42 percentage points and \Suseful\ by 10.75 points relative to central-only unlimited routing, but increase mean transferred data by approximately 1,438 Mbit per case. The comparison demonstrates that link availability and central-node fraction must be interpreted jointly.

\section{Task-Size Sensitivity}

Task-size scaling provides a direct communication-load stress test while retaining the same tasks and physical opportunities. Table~\ref{tab:task_size} and Figure~\ref{fig:task_size} show monotonic degradation as packets grow. Halving the packet size slightly improves performance and approximately halves traffic; quadrupling it reduces observation fulfilment by 8.14 points and strict completion by 14.59 points relative to the 1$\times$ case.

\begin{table}[htbp]
\centering
\caption{California task-size sensitivity over 890 complete architectures.}
\label{tab:task_size}
\begin{tabular}{rrrrr}
\toprule
Multiplier & Obs. (\%) & \Suseful\ (\%) & Coverage (\%) & Data (Mbit/case) \\
\midrule
0.5$\times$ & 79.33 & 76.01 & 91.72 & 1,036.55 \\
1$\times$ & 77.62 & 74.69 & 91.14 & 2,047.57 \\
2$\times$ & 75.12 & 71.73 & 89.38 & 3,979.20 \\
4$\times$ & 69.48 & 60.10 & 82.36 & 7,219.41 \\
\bottomrule
\end{tabular}
\end{table}

\begin{figure}[htbp]
\centering
\includegraphics[width=\textwidth]{05_task_size_sensitivity.png}
\caption{Task-size sensitivity of performance and communication burden.}
\label{fig:task_size}
\end{figure}

The balanced 40-percent-CN candidate retains 96.875 percent observation at 2$\times$ size but falls to 87.5 percent at 4$\times$; its \Suseful\ changes from 100 percent at 1$\times$ to 96.875 and 78.125 percent. A recommendation based on small task descriptions is therefore not invariant to heavier products.

\section{Priority-Class Outcomes}

California contains one Critical, six High, sixteen Medium, and nine Low tasks. Figure~\ref{fig:priority_outcomes} shows that urgent classes are substantially harder under both anchor scenarios. The nominal mean observation outcomes are 8.20, 23.09, 63.43, and 89.90 percent for Critical through Low. Under the unlimited replay they become 29.21, 53.75, 79.66, and 95.31 percent. Corresponding unlimited \Suseful\ means are 24.49, 54.55, 77.80, and 88.16 percent.

\begin{figure}[htbp]
\centering
\includegraphics[width=0.92\textwidth]{07_priority_outcomes.png}
\caption{California task-priority outcomes under nominal and unlimited scenarios.}
\label{fig:priority_outcomes}
\end{figure}

This pattern does not isolate a causal priority effect. Priority is jointly encoded through deadline and message size, and each class contains different task locations and release times. The Critical class has one task. The result supports a workload-heterogeneity conclusion, not a general claim that high-FRP fires are intrinsically harder to observe.

\section{Failure Robustness}

For failure scenarios, each architecture is compared with its own unlimited baseline. Table~\ref{tab:robustness_30} reports the average percentage-point effect at the 30-percent level. Ground-station outage is the dominant tested stressor. Satellite and central-node failure produce moderate average losses; capacity degradation is smaller at the tested level.

\begin{table}[htbp]
\centering
\caption{Mean paired performance change at the 30-percent stress level.}
\label{tab:robustness_30}
\begin{tabularx}{\textwidth}{L{0.40\textwidth}rrY}
\toprule
Failure family & $\Delta$ Obs. (pp) & $\Delta$ \Suseful\ (pp) & Interpretation \\
\midrule
Ground-station outage & $-21.46$ & $-20.84$ & Dominant tested vulnerability \\
Targeted central-node failure & $-2.95$ & $-6.57$ & Larger dissemination than observation penalty \\
Random central-node failure & $-2.86$ & $-5.98$ & Moderate mean loss \\
Random satellite failure & $-3.01$ & $-3.91$ & Geometry and forwarding loss \\
Capacity degradation & $< -1$ & approximately $-1$ & Small at tested severity \\
\bottomrule
\end{tabularx}
\end{table}

\begin{figure}[htbp]
\centering
\includegraphics[width=\textwidth]{06_robustness_delta_s100.png}
\caption{Paired change in strict useful dissemination across Fresh V2 robustness families and severities.}
\label{fig:robustness_s100}
\end{figure}

For the balanced candidate \texttt{T060\_P10\_H0800\_I0986\_CN040\_F1}, the unlimited baseline is 96.875 percent observation and 100 percent \Suseful. At 30-percent random central-node failure, replicate means are 91.25 and 93.75 percent; at 30-percent random satellite failure they are 92.08 and 95.52 percent. Deterministic targeted-CN failure retains 93.75 percent observation and 100 percent \Suseful, although P100 $T_{100}$ increases from 66.10 to 124.42 minutes. Ground-station-outage replicate means fall to 66.77 and 69.79 percent. Ground diversity is therefore a more direct design need than simply raising the central-node fraction beyond the intermediate range.

\section{Pareto Trade Space}

No architecture dominates all others on performance and burden. Figure~\ref{fig:v2_pareto} shows the unlimited trade space. The original six-objective analysis---observation, \Suseful, useful coverage, satellites, central nodes, and transferred data---contains 48 exact Pareto cases. Adding P100 $T_{100}$ and transfer-event count for consistency with the utility definition produces an eight-objective frontier of 108 cases; all original 48 remain efficient.

\begin{figure}[htbp]
\centering
\includegraphics[width=\textwidth]{09_unlimited_pareto_trade_space.png}
\caption{California unlimited performance--resource trade space and exact Pareto membership.}
\label{fig:v2_pareto}
\end{figure}

Table~\ref{tab:decision_cases} gives three transparent decision points. The balanced candidate achieves complete useful dissemination and coverage with 60 satellites and less than one-sixth of the transferred data of the fastest perfect-performance extreme.

\begin{table}[htbp]
\centering
\small
\caption{Illustrative California Fresh V2 Pareto candidates.}
\label{tab:decision_cases}
\begin{tabularx}{\textwidth}{L{0.31\textwidth}rrrrr}
\toprule
Case & Obs. (\%) & \Suseful\ (\%) & Cov. (\%) & P100 $T_{100}$ (min) & Mbit \\
\midrule
\shortstack[l]{\texttt{T060\_P10\_H0800}\protect\\\texttt{I0986\_CN040\_F1}} & 96.875 & 100 & 100 & 66.10 & 888.5 \\
\shortstack[l]{\texttt{T060\_P05\_H0800}\protect\\\texttt{I0800\_CN070\_F1}} & 100 & 100 & 100 & 56.20 & 1,422 \\
\shortstack[l]{\texttt{T180\_P10\_H0800}\protect\\\texttt{I0800\_CN100\_F1}} & 100 & 100 & 100 & 33.59 & 5,670 \\
\bottomrule
\end{tabularx}
\end{table}

The first case has 24 central nodes; the second has 42; the third has 180. The first is the recommended California compromise under the declared balanced weighting, not a universal optimum or a cost-optimal spacecraft design.

\section{Objective-Weight Analysis}

The weight analysis introduces preference only after exact Pareto screening and normalization over all 890 complete cases. Table~\ref{tab:weight_profiles} shows three declared high-level profiles. Cost-dominant is shorthand for resource-dominant; the score is non-monetary.

\begin{table}[htbp]
\centering
\caption{Declared weight profiles and selected California architecture.}
\label{tab:weight_profiles}
\begin{tabularx}{\textwidth}{L{0.20\textwidth}rrY}
\toprule
Profile & Performance & Resource & Winner \\
\midrule
Resource-dominant & 0.25 & 0.75 & \texttt{T060\_P05\_H1000\_I0986\_CN030\_F1} \\
Balanced & 0.50 & 0.50 & \texttt{T060\_P10\_H0800\_I0986\_CN040\_F1} \\
Performance-dominant & 0.75 & 0.25 & \texttt{T060\_P10\_H0800\_I0986\_CN040\_F1} \\
\bottomrule
\end{tabularx}
\end{table}

\begin{figure}[htbp]
\centering
\includegraphics[width=\textwidth]{objective_weight_sweep.png}
\caption{Winner transitions as the high-level performance weight $\lambda$ changes from zero to one.}
\label{fig:objective_sweep}
\end{figure}

The 40-percent candidate wins for $0.31\le\lambda\le0.78$. Above 0.78 the winner moves first to a 70-percent-CN case and then to 100-percent-CN cases, reaching the 180-satellite extreme only near pure performance weighting. Pure resource weighting selects a CN0 case with no observation or strict-dissemination success, which is operationally unacceptable. The endpoint shows why feasibility thresholds must precede utility ranking.

\begin{figure}[htbp]
\centering
\includegraphics[width=0.92\textwidth]{objective_weight_winner_frequency.png}
\caption{Winner frequency across 10,000 randomized objective-weight draws (seed 20260825).}
\label{fig:objective_mc}
\end{figure}

With randomized high-level and internal weights, the balanced 40-percent case wins 41.79 percent of 10,000 draws. Other frequent winners include a 60-satellite 100-percent-CN case (12.11 percent), a 20-percent-CN case (11.49 percent), a CN0 case (8.40 percent), and the 70-percent perfect-performance case (7.91 percent). This supports a broad intermediate solution under the declared preference distribution, while also showing that hard requirements and stakeholder weights can change the choice.

\section{California Result Synthesis}

Fresh V2 supports five connected conclusions. First, the nominal copy/hop limits are a major dissemination bottleneck. Second, an intermediate CN fraction is preferable under the bounded policy, while unlimited performance saturates as communication burden continues to grow. Third, ground-station outage is the dominant tested failure. Fourth, message size materially changes the attainable solution. Fifth, the 40-percent-CN, 60-satellite candidate is a strong California compromise across a wide weight interval, but it is conditional on the implemented task set, fixed observation radius, Munich ground station, and non-monetary resource model.

The open Fresh V2 evidence is equally important: India has not yet passed the same analysis, observation-radius sensitivity has not been run, and separate first-reception latency is not present in the canonical case table. These limitations define the remaining work rather than weakening the validated California claims.
